\begin{lemma}
\label{lemma-quasi-coherence-direct-image}
Let $f : X \to S$ be a morphism of schemes.
Assume that $f$ is quasi-separated and quasi-compact.
\begin{enumerate}
\item The functor $Rf_*$ sends $D_\QCoh(\mathcal{O}_X)$
into $D_\QCoh(\mathcal{O}_S)$.
\item If $S$ is quasi-compact, there exists an integer $N = N(X, S, f)$
such that for an object $E$ of $D_\QCoh(\mathcal{O}_X)$
with $H^m(E) = 0$ for $m > 0$ we have
$H^m(Rf_*E) = 0$ for $m \geq N$.
\item In fact, if $S$ is quasi-compact we can find $N = N(X, S, f)$
such that for every morphism of schemes $S' \to S$
the same conclusion holds for the functor $R(f')_*$
where $f' : X' \to S'$ is the base change of $f$.
\end{enumerate}
\end{lemma}

\begin{proof}
Let $E$ be an object of $D_\QCoh(\mathcal{O}_X)$. To prove (1) we have to
show that $Rf_*E$ has quasi-coherent cohomology sheaves. The question is local
on $S$, hence we may assume $S$ is quasi-compact. Pick $N = N(X, S, f)$ as in
Cohomology of Schemes, Lemma
\ref{coherent-lemma-quasi-coherence-higher-direct-images}.
Thus $R^pf_*\mathcal{F} = 0$ for all quasi-coherent $\mathcal{O}_X$-modules
$\mathcal{F}$ and all $p \geq N$ and the same remains true after base change.

\medskip\noindent
First, assume $E$ is bounded below. We will show (1) and (2) and (3) hold
for such $E$ with our choice of $N$. In this case we can for example use the
spectral sequence
$$
R^pf_*H^q(E) \Rightarrow R^{p + q}f_*E
$$
(Derived Categories, Lemma \ref{derived-lemma-two-ss-complex-functor}),
the quasi-coherence of $R^pf_*H^q(E)$, and the vanishing of $R^pf_*H^q(E)$
for $p \geq N$ to see that (1), (2), and (3) hold in this case.

\medskip\noindent
Next we prove (2) and (3). Say $H^m(E) = 0$ for $m > 0$.
Let $U \subset S$ be affine open. By Cohomology of Schemes, Lemma
\ref{coherent-lemma-quasi-coherence-higher-direct-images-application}
and our choice of $N$
we have $H^p(f^{-1}(U), \mathcal{F}) = 0$ for $p \geq N$
and any quasi-coherent $\mathcal{O}_X$-module $\mathcal{F}$.
Hence we may apply Lemma \ref{lemma-application-nice-K-injective}
to the functor $\Gamma(f^{-1}(U), -)$ to see that
$$
R\Gamma(U, Rf_*E) = R\Gamma(f^{-1}(U), E)
$$
has vanishing cohomology in degrees $\geq N$. Since this holds for
all $U \subset S$ affine open we conclude that $H^m(Rf_*E) = 0$
for $m \geq N$.

\medskip\noindent
Next, we prove (1) in the general case. Recall that there is a
distinguished triangle
$$
\tau_{\leq -n - 1}E \to E \to \tau_{\geq -n}E \to
(\tau_{\leq -n - 1}E)[1]
$$
in $D(\mathcal{O}_X)$, see Derived Categories, Remark
\ref{derived-remark-truncation-distinguished-triangle}.
By (2) we see that $Rf_*\tau_{\leq -n - 1}E$
has vanishing cohomology sheaves in degrees $\geq -n + N$.
Thus, given an integer $q$ we see that $R^qf_*E$ is equal
to $R^qf_*\tau_{\geq -n}E$ for some $n$ and the result
above applies.
\end{proof}
